\section{Câu 2. Giải bài toán từng bước}

Ta dùng lan truyền ràng buộc (Forward Checking) sau mỗi phép gán. Khi việc
lan truyền chưa xác định được nghiệm duy nhất, ta phân nhánh theo biến có
miền giá trị nhỏ và kiểm tra các ràng buộc liên quan.

\subsection{Bước 1: Xác định vị trí của Alice}

Từ ràng buộc $C_1=((A),\{1\})$, ta có
\[
    D_A=\{1\},\qquad A=1.
\]
Áp dụng nhóm ràng buộc AllDiff, loại ghế 1 khỏi miền của các biến còn lại:
\[
    D_B=D_C=D_D=D_E=D_F=\{2,3,4,5,6\}.
\]

\subsection{Bước 2: Xác định vị trí của Bob}

Theo ràng buộc đối diện $C_2$, khi $A=1$ thì chỉ có cặp $(1,4)$ thỏa mãn.
Do đó
\[
    D_B=\{4\},\qquad B=4.
\]
Tiếp tục áp dụng AllDiff, ta được
\[
    D_C=D_D=D_E=D_F=\{2,3,5,6\}.
\]

\subsection{Bước 3: Thu hẹp miền của Frank và Dave}

Vì Alice ở ghế 1, hai ghế kề Alice là 2 và 6. Từ $C_3$, Frank không thể
ngồi ở hai ghế này:
\[
    D_F=\{3,5\}.
\]

Vì Bob ở ghế 4, Dave phải ngồi cạnh ghế 4. Hai lựa chọn ban đầu là ghế 3
và 5, do đó
\[
    D_D=\{3,5\}.
\]

\subsection{Bước 4: Thu hẹp miền của Eve và Carol}

Hai biến $D$ và $F$ cùng có miền $\{3,5\}$, đồng thời $C_{19}$ yêu cầu
$D\ne F$. Vì vậy hai ghế 3 và 5 chắc chắn được phân cho Dave và Frank.
Loại chúng khỏi miền của Eve và Carol, ta thu được
\[
    D_E=D_C=\{2,6\}.
\]

\subsection{Bước 5: Phân nhánh và tìm nghiệm}

Chọn biến Dave vì $D_D=\{3,5\}$ và Dave liên kết với nhiều ràng buộc.

\subsubsection{Trường hợp 1: $D=3$}

Do $D\ne F$ và $D_F=\{3,5\}$, suy ra $F=5$. Dave phải ngồi cạnh Eve;
các ghế kề 3 là 2 và 4, nhưng ghế 4 đã thuộc về Bob, nên $E=2$. Cuối
cùng, từ AllDiff và $D_C=\{2,6\}$, suy ra $C=6$.

Nghiệm chính là
\[
    (A,B,C,D,E,F)=(1,4,6,3,2,5).
\]

\begin{table}[htbp]
    \centering
    \caption{Nghiệm chính của bài toán CSP}
    \label{tab:nghiem-csp}
    \renewcommand{\arraystretch}{1.25}
    \begin{tabular}{c|cccccc}
        \toprule
        Nhân vật / Ghế & 1 & 2 & 3 & 4 & 5 & 6 \\
        \midrule
        Alice & $\checkmark$ &  &  &  &  &  \\
        Bob   &  &  &  & $\checkmark$ &  &  \\
        Carol &  &  &  &  &  & $\checkmark$ \\
        Dave  &  &  & $\checkmark$ &  &  &  \\
        Eve   &  & $\checkmark$ &  &  &  &  \\
        Frank &  &  &  &  & $\checkmark$ &  \\
        \bottomrule
    \end{tabular}
\end{table}

\subsubsection{Trường hợp 2: $D=5$}

Đây là nghiệm đối xứng gương của trường hợp trên. Khi đó
\[
    (A,B,C,D,E,F)=(1,4,2,5,6,3).
\]

\subsection{Kết luận}

Bài toán có hai nghiệm đối xứng nếu chỉ xét các ràng buộc đã cho. Nghiệm
chính được trình bày trong Bảng~\ref{tab:nghiem-csp}: Alice ngồi ghế 1,
Eve ghế 2, Dave ghế 3, Bob ghế 4, Frank ghế 5 và Carol ghế 6. Nghiệm còn
lại là ảnh đối xứng qua trục đi qua ghế 1 và ghế 4.

